\documentclass[12pt]{article}

\usepackage[left=2.5cm,top=2.5cm,right=2.5cm,bottom=2.5cm]{geometry}

%% Macros and packages go here

\begin{document}
\title{COMP0005 Algorithms: Group Coursework Report}
\author{Group Number: XXX}
\date{Academic year 2025/26}
\maketitle



\section{Experimental framework overview}
\emph{(Section 1 should be about one page in length.)}

\subsection{Framework design and architecture}

The framework tests the insertion and search time as the number of nodes in a tree grows.
The synthetic keys are generated in the $TestDataGenerator$.
To suitably stress test the trees, four different methods of generating keys are used: sorted, reverse sorted, uniform random, and normal random. 
Uniform random data is the best case, where keys are naturally self balancing [REF]. 
Normal random data is the realistic average case, as due to Central Limit Theorem [REF] most natural data takes a normal distribution.
Sorted and reverse sorted are the worst case, where the cost of rebalancing is most expressed [REF] due to always inserting on the very left or right of the tree.

As the time to compare keys is a control variable, the framework ensures all keys are a constant string length.
This makes the time to compare keys $O(k)$, where k is the length of the string. 
As $k$ is constant, the cost of key comparisons is $O(1)$ with respect to the number of nodes in the tree, ensuring the theoretical relationship is still $O(logn)$.
Strings are compared lexicographically, not numerically, so casting to string would change the intended distribution of random data. To maintain numerical ordering, truncating and left padding with 0's is used.
The string length of 7 was chosen as it allows for $10^7$ unique keys, more than the number of nodes being inserted.

The data is collected using the $ExperimentalFramework$. The algorithm to time inserts and searches simultaneously builds the tree while timing operations at fixed intervals of size $i$. 
Operations are timed singularly, not cumulatively. At every interval, $k$ inserts and searches are timed for statistical significance, 
with the key being deleted from the tree after every insert (not timed), so that the tree returns to the size being measured. 
A search for a random key in, and not in the tree is performed, to measure both the average and worst cases to search for a key.
The tree is built until it reaches a maximum size $n$, with $n$ being sufficiently large to measure the relationship between tree size and operation time.
This process is repeated $m$ times, to minimize bias from OS processes such as scheduling.

After collecting the data, it is processed using the $StatisticsFramework$.
The framework calculates the log values of the tree's size as the theoretical $O(logn)$ complexity for insert and search would appear as a linear relation on a log plot.
This allows the framework to calculate a value for PMCC, a quantitative measurement of linearity.
The regression line is calculated so that its gradient can compare the hidden constants of different tree implementations.
The RMSE is calculated to measure how close data points are to the regression line. The allows for reasoning about the predictability of each implementations performance.
The residuals are plotted as it can reveal behaviours about a tree, 
for example oscillations of residuals may indicate rebalancing at certain sizes, and a lack of such at other sizes.

\subsection{Hardware/software setup for experiments}




\section{Performance results}
\emph{(Section 2 should be about two to two-and-a-half pages in length.)}

\subsection{[First chosen data structure]}

\subsection{[Second chosen data structure]}

\subsection{[Third chosen data structure]}




\section{Comparative assessment}
\emph{(Section 3 should be about one to one-and-a-half pages in length.)}



\section{Team contributions}
\emph{(Section 4 should be no more than half a page in length. Express work share as a percentage totalling to $100$: for example, four students who all contributed equally should receive $25\%$.)}


\vspace{1cm}

\begin{tabular}{|l|l|l|l|}
\hline
{\bf Name} & {\bf Portico ID} & {\bf Key contributions} & {\bf Work share (\%)} \\[1mm]
\hline
(name)    &  (number)  & (text) & ($xx\%$) \\\hline
          &            &         &          \\
%% more entries as needed
\hline
\end{tabular}


\end{document}
